%======================================================================
\chapter{Correspondence with the formal proof}\label{app:formal}
%======================================================================

The table lists the definitions and results of this volume that carry a ``Formal counterpart''
line, together with the corresponding declarations of the Lean development. All declaration
names start with \lean{HPC.}, except the Mathlib names marked as such. The formal statements are
sometimes slightly more general, or are split differently into lemmas; the table points to the
place where the content of each result is proved. Rows are in the order of the text.

The main theorems of the volume are, in the Lean development (prefix \lean{HPC.} omitted):
\begin{center}
\begin{tabular}{@{}p{0.42\textwidth}p{0.54\textwidth}@{}}
\toprule
Unicoherence of the plane (Theorem~\ref{thm:unicoherent}) & \lean{Moore.plane_unicoherent}\\
Brouwer's theorem (Theorem~\ref{thm:brouwer}) & \lean{Moore.brouwer_boundary}\\
Janiszewski's theorem (Theorem~\ref{thm:janiszewski}) & \lean{Moore.janiszewski_compact}\\
Borsuk's theorem (Theorem~\ref{thm:borsuk}) & \lean{Moore.borsuk_exists_log}\\
Invariance of plane separation (Theorem~\ref{thm:sep-invariant}) & \lean{Jones.separatesPlane_of_homeomorph}\\
Jordan curve theorem for polygons (Theorem~\ref{thm:jct-poly}) & \lean{Plane.jct_polygon}\\
Enclosing polygons (Theorem~\ref{thm:enclose}) & \lean{Plane.enclose_holds}\\
Tree-like continua contain no separating continuum (Theorem~\ref{thm:treelike-nonsep}) &
\lean{Plane.treeLike_nonseparating}\\
Logarithms on tree-like continua (Theorem~\ref{thm:log-treelike}) & \lean{Hagopian.exists_log_of_isTreeLike}\\
Bing's Theorem~6 (Theorem~\ref{thm:bing6}) & \lean{Plane.theorem6_direct}, \lean{Plane.theorem6_1951}\\
Bing's characterization (Corollary~\ref{cor:treelike-iff}) & \lean{Plane.isTreeLike_iff_nonseparating}\\
Weak invariance (Theorem~\ref{thm:weak-invariance}) &
\lean{Hagopian.exists_separating_subcontinuum_of_homeomorph}\\
Bing's Theorem~7 (Theorem~\ref{thm:bing7}) & \lean{Plane.theorem7_1951}\\
\bottomrule
\end{tabular}
\end{center}
All of them are proved in the Lean development without \lean{sorry} and without axioms beyond the
standard ones of Lean and Mathlib. The direct proof of Theorem~\ref{thm:bing6}
(\lean{Plane.theorem6_direct} and the lemmas \lean{Plane.janiszewski_polygonal_of_compact},
\lean{Plane.exists_theta_point_open}, \lean{Plane.compFamily_no_circular_open},
\lean{Plane.treeChain_of_cover_open}) is in the file \texttt{HPC/Plane/Theorem6Direct.lean}, which was
written for these notes.

\begin{longtable}{@{}p{0.40\textwidth}p{0.56\textwidth}@{}}
\toprule
In these notes & In the Lean development\\
\midrule
\endhead
\bottomrule
\endfoot
\multicolumn{2}{@{}l}{\emph{Chapter~\ref{ch:log}}}\\
Lemma~\ref{lem:int-valued} & \lean{HPC.Plane.eq_of_isPreconnected_int_valued}.\\
Uniqueness (Lemma~\ref{lem:log-unique}) & \lean{HPC.Plane.sub_const_of_exp_eq}.\\
Principal logarithm (Section~\ref{sec:log-unique}) & Mathlib's \lean{Complex.log}, \lean{Complex.arg}; the second fact is \lean{HPC.Hagopian.log_mul_near_one}.\\
Nearby functions (Lemma~\ref{lem:log-near}) & \lean{HPC.Moore.exists_log_of_near}.\\
Homotopy lifting; quoted (Theorem~\ref{thm:lifting}) & Mathlib's \lean{IsCoveringMap.liftHomotopy} applied to \lean{Complex.isCoveringMap_exp}.\\
Logarithms on contractible spaces (Theorem~\ref{thm:contractible}) & \lean{HPC.Plane.exists_exp_lift}.\\
Corollary~\ref{cor:log-exist} & (a) \lean{HPC.Moore.exists_log_pt}; (b) \lean{HPC.Moore.exists_log_of_convex}; (c) \lean{HPC.Moore.exists_exp_lift_of_homotopy}.\\
A logarithm of $z-p$ on a small ball (Lemma~\ref{lem:log-ball}) & \lean{HPC.Plane.exists_log_on_ball}.\\
Unicoherence of the plane (Theorem~\ref{thm:unicoherent}) & \lean{HPC.Moore.plane_unicoherent}; the same argument for a compact space contractible by an explicit homotopy is \lean{HPC.Plane.isUnicoherent_univ_of_contraction}.\\
Lemma~\ref{lem:frontier-component} & \lean{HPC.Moore.frontier_connectedComponentIn_subset}, \lean{HPC.Plane.isOpen_connectedComponentIn_of_isOpen}.\\
Brouwer (Theorem~\ref{thm:brouwer}) & \lean{HPC.Moore.brouwer_boundary}.\\[4pt]
\multicolumn{2}{@{}l}{\emph{Chapter~\ref{ch:sep}}}\\
Definition~\ref{def:separates} & \lean{HPC.Plane.SeparatesPlane}, \lean{HPC.Plane.SeparatesPts} (\lean{connectedComponentIn} for $\comp Ax$).\\
Remarks after Definition~\ref{def:separates} & \lean{IsOpen.isConnected_iff_isPathConnected} (Mathlib); \lean{HPC.Plane.separatesPts_mono}.\\
Lemma~\ref{lem:compl-square} & \lean{HPC.Plane.isConnected_compl_closedBall_pt}, \lean{HPC.Plane.isPreconnected_compl_closedBall}.\\
Lemma~\ref{lem:unbounded-comp} & (a) \lean{HPC.Moore.mem_component_of_unbounded}; (b) \lean{HPC.Plane.exists_bounded_component_of_separates}.\\
Same component $\Rightarrow$ logarithm (Theorem~\ref{thm:log-same-comp}) & \lean{HPC.Moore.exists_log_of_mem_component}.\\
Bounded component $\Rightarrow$ no logarithm (Theorem~\ref{thm:nolog-bounded}) & \lean{HPC.Moore.no_log_pow_of_bounded_component}.\\
Corollary~\ref{cor:nolog-bounded} & \lean{HPC.Moore.no_log_of_bounded_component} (proved directly, by the same argument).\\
Janiszewski (Theorem~\ref{thm:janiszewski}) & \lean{HPC.Moore.janiszewski_compact}.\\
Janiszewski for polygonal sets (Corollary~\ref{cor:janiszewski-poly}) & \lean{HPC.Plane.janiszewski_polygonal_of_compact}. The formal development also contains an independent combinatorial proof, \lean{HPC.Plane.janiszewski_polygonal}, by mod-2 cycles of straight-line graphs; see Remark~\ref{rem:jan-cycles}.\\
Definition~\ref{def:good-zero} & \lean{HPC.Moore.IsGoodZero}, \lean{HPC.Moore.sgnConj}.\\
Lemma~\ref{lem:good-zero} & (a) \lean{HPC.Moore.IsGoodZero.isolated}; (b) \lean{HPC.Moore.isGoodZero_of_hasFDerivAt}; (c) \lean{HPC.Moore.IsGoodZero.mul}; (d) \lean{HPC.Moore.finite_zeros_of_goodZeros}.\\
Removing a good zero (Lemma~\ref{lem:remove-zero}) & \lean{HPC.Moore.IsGoodZero.remove}.\\
Lemma~\ref{lem:log-factor} & \lean{HPC.Moore.exists_log_sgnConj}.\\
Induction on the zeros (Lemma~\ref{lem:zeros-induction}) & \lean{HPC.Moore.exists_log_of_goodZeros}.\\
Stone--Weierstrass and Sard (quoted) & Mathlib's \lean{ContinuousMap.exists_mem_subalgebra_near_continuous_of_isCompact_of_separatesPoints} and \lean{MeasureTheory.addHaar_image_eq_zero_of_det_fderivWithin_eq_zero}; in the development \lean{HPC.Moore.exists_contDiff_near} and \lean{HPC.Moore.exists_small_regular_value}.\\
Borsuk (Theorem~\ref{thm:borsuk}) & \lean{HPC.Moore.borsuk_exists_log}.\\
Theorem~\ref{thm:sep-invariant} & \lean{HPC.Jones.separatesPlane_of_homeomorph}.\\[4pt]
\multicolumn{2}{@{}l}{\emph{Chapter~\ref{ch:poly}}}\\
Definition~\ref{def:polygon} & \lean{HPC.Plane.polygonOf}, \lean{HPC.Plane.polyEdge}, \lean{HPC.Plane.IsSimpleVertices}, \lean{HPC.Plane.IsSimpleClosedPolygon} (with \lean{HPC.Plane.isSimpleClosedPolygon_iff}).\\
Definition~\ref{def:ins} & \lean{HPC.Plane.polyInterior}.\\
Lemma~\ref{lem:ins-basic} & (a) \lean{HPC.Plane.polyInterior_disjoint_of_unbounded}; (b) \lean{HPC.Plane.polyInterior_eq_of_split} (with \lean{HPC.Plane.connectedComponentIn_eq_of_split}).\\
The half-open crossing rule (Definition~\ref{def:crosses}) & \lean{HPC.Plane.CrossesEast}, \lean{HPC.Plane.crossCount}, \lean{HPC.Plane.parity}.\\
Vertex jumps & \lean{HPC.Plane.vJump}, \lean{HPC.Plane.ceInd}.\\
The vertex rule (Lemma~\ref{lem:vertex-rule}) & \lean{HPC.Plane.crossesEast_jump} (with \lean{HPC.Plane.crossesEast_jump_end}).\\
Local constancy of the parity (Theorem~\ref{thm:parity-loc-const}) & \lean{HPC.Plane.parity_eventually_eq}, \lean{HPC.Plane.parity_locallyConstant}, \lean{HPC.Plane.sum_vJump_telescope}, \lean{HPC.Plane.parity_eq_of_isPreconnected}, \lean{HPC.Plane.isOpen_parity_set}, \lean{HPC.Plane.parity_far}, \lean{HPC.Plane.parity_eq_zero_of_far}.\\
Lemma~\ref{lem:local-rays} & \lean{HPC.Plane.polygon_local_rays}, \lean{HPC.Plane.rayFrom}.\\
Two rays cut a square into two pieces (Lemma~\ref{lem:ball-rays}) & \lean{HPC.Plane.ball_diff_rays}.\\
Local sides from a splitting (Lemma~\ref{lem:local-side}) & \lean{HPC.Plane.local_side_of_split}.\\
Gluing local pieces (Lemma~\ref{lem:local-meet}) & \lean{HPC.Plane.isPreconnected_biUnion_of_local_meet}.\\
Flips (Lemma~\ref{lem:flip}) & (a) \lean{HPC.Plane.parity_flip_at} (and \lean{HPC.Plane.exists_parity_flip}); (b) \lean{HPC.Plane.parity_flip_vertical}.\\
Proposition~\ref{prop:both-parities} & \lean{HPC.Plane.mem_closure_parity_set} (with \lean{HPC.Plane.exists_near_on_segment}).\\
First hit (Lemma~\ref{lem:first-hit}) & \lean{HPC.Plane.exists_first_hit}.\\
Jordan curve theorem for simple closed polygons (Theorem~\ref{thm:jct-poly}) & \lean{HPC.Plane.jct_polygon}; the steps are \lean{HPC.Plane.polygon_separates}, \lean{HPC.Plane.frontier_parity_set}, \lean{HPC.Plane.local_piece}, \lean{HPC.Plane.sideNbhd}, \lean{HPC.Plane.sideNbhd_isConnected}, \lean{HPC.Plane.exists_segment_to_sideNbhd}, \lean{HPC.Plane.isConnected_parity_set}.\\
Corollary~\ref{cor:jct-consequences} & (a) \lean{HPC.Plane.polyInterior_eq_oddSet}, \lean{HPC.Plane.PlaneFacts.polyInterior_eq}; (b) \lean{HPC.Plane.PlaneFacts.compl_polyInterior}; (c) \lean{HPC.Plane.PlaneFacts.local_side}; (d) \lean{HPC.Plane.sides_of_split} and \lean{HPC.Plane.PlaneFacts.sides}. The bundle of plane facts \lean{HPC.Plane.PlaneFacts} is proved as \lean{HPC.Plane.planeFacts_holds}.\\
Definition~\ref{def:polypath} & \lean{HPC.Plane.polyPathOf}, \lean{HPC.Plane.IsSimplePolyArc} (and the vertex-function form \lean{HPC.Plane.vpath}, \lean{HPC.Plane.IsSArc}).\\
Proposition~\ref{prop:polypath} & (a) \lean{HPC.Plane.exists_vpath_of_isConnected}, \lean{HPC.Plane.exists_polyPath_of_isConnected}; (b) \lean{HPC.Plane.exists_simplePolyArc_of_polyPath} (\lean{HPC.Plane.exists_sarc_of_vpath}); (c) \lean{HPC.Plane.IsSimplePolyArc.exists_first_mem}.\\
Closed chains of arcs (Lemma~\ref{lem:closed-chain}) & \lean{HPC.Plane.closed_chain_polygon}; for three arcs, the threaded form \lean{HPC.Plane.IsThreaded.isSimpleClosedPolygon}.\\
Proposition~\ref{prop:open-minus-point} & \lean{HPC.Plane.isConnected_diff_singleton_of_isOpen} (with \lean{HPC.Plane.isPreconnected_ball_diff_singleton}).\\
Theorem~\ref{thm:nolog-polygon} & \lean{HPC.Plane.no_log_on_polygon}, from \lean{HPC.Plane.log_diff_edge}, \lean{HPC.Plane.arg_div_jump}, \lean{HPC.Plane.crossesEast_iff_cross} and \lean{HPC.Plane.edge_jump}.\\[4pt]
\multicolumn{2}{@{}l}{\emph{Chapter~\ref{ch:enc}}}\\
Definition~\ref{def:graph} & \lean{HPC.Plane.Seg}, \lean{HPC.Plane.graphSet}, \lean{HPC.Plane.vertices}, \lean{HPC.Plane.deg}, \lean{HPC.Plane.IsProper}, \lean{HPC.Plane.IsCycle}.\\
Parity of a graph & \lean{HPC.Plane.cycParity}.\\
Additivity of the parity & \lean{HPC.Plane.cycParity_symmDiff}, \lean{HPC.Plane.cycParity_foldr}, \lean{HPC.Plane.cycParity_far}, \lean{HPC.Plane.cycParity_polyEdges}.\\
Proposition~\ref{prop:cycle-parity} & \lean{HPC.Plane.sum_vJump_even}, \lean{HPC.Plane.cycParity_locallyConstant}, \lean{HPC.Plane.cycParity_eq_of_isPreconnected}, \lean{HPC.Plane.separatesPts_of_cycParity_ne}; the necessity of nondegeneracy is \lean{HPC.Plane.not_sum_vJump_even_skeleton}.\\
Flip across an edge of a cycle (Lemma~\ref{lem:cycle-flip}) & \lean{HPC.Plane.cycParity_flip} (from \lean{HPC.Plane.cycParity_flip_horizontal_shift} and \lean{HPC.Plane.cycParity_flip_vertical_shift}).\\
Lemma~\ref{lem:side-strip} & \lean{HPC.Plane.sideRadius}, \lean{HPC.Plane.edge_local}, \lean{HPC.Plane.sideStrip}, \lean{HPC.Plane.sideStrip_isPreconnected}, \lean{HPC.Plane.sideStrip_subset_compl}, \lean{HPC.Plane.subset_or_disjoint_component}.\\
Odd set & \lean{HPC.Plane.oddSet}.\\
Lemma~\ref{lem:odd-even} & \lean{HPC.Plane.even_deg_oddSet}, \lean{HPC.Plane.isCycle_oddSet}; the circle argument is \lean{HPC.Plane.vertex_arcs} and \lean{HPC.Plane.zmod2_jump_sum}.\\
Theorem~\ref{thm:odd-cycle} & \lean{HPC.Plane.gFun}, \lean{HPC.Plane.gFun_eq_across_edge}, \lean{HPC.Plane.gFun_const}, \lean{HPC.Plane.cycParity_oddSet_ne}.\\
Proposition~\ref{prop:simple-cycles} & \lean{HPC.Plane.exists_cycle_vertices}, \lean{HPC.Plane.simpleCycle_of_vertices}, \lean{HPC.Plane.exists_simpleCycle_decomposition}, \lean{HPC.Plane.IsSimpleCycle}, \lean{HPC.Plane.polyInterior_eq_oddSet}.\\
Lemma~\ref{lem:grid} & \lean{HPC.Plane.gsq}, \lean{HPC.Plane.sqUnion}, \lean{HPC.Plane.bdryOf}, \lean{HPC.Plane.isProper_bdryOf}, \lean{HPC.Plane.mem_frontier_sqUnion_iff}, \lean{HPC.Plane.mem_graphSet_bdryOf_iff}, \lean{HPC.Plane.frontier_sqUnion}, \lean{HPC.Plane.mem_interior_sqUnion}.\\
Lemma~\ref{lem:grid-union} & \lean{HPC.Plane.sqIdx}, \lean{HPC.Plane.gridUnion}, \lean{HPC.Plane.bdryGraph}, \lean{HPC.Plane.subset_interior_gridUnion}, \lean{HPC.Plane.disjoint_frontier_gridUnion}, \lean{HPC.Plane.gridUnion_subset_cthickening}, \lean{HPC.Plane.frontier_gridUnion}, \lean{HPC.Plane.farPt}, \lean{HPC.Plane.halfPlane_subset_component}, \lean{HPC.Plane.not_isBounded_component}, \lean{HPC.Plane.interior_gridUnion_disjoint_component}.\\
Theorem~\ref{thm:enclosing-polygon} & \lean{HPC.Plane.cycParity_oddSet_eq_one_on}, \lean{HPC.Plane.exists_enclosing_polygon}, \lean{HPC.Plane.exists_enclosing_polygon_pt}.\\
Enclosing polygons (Theorem~\ref{thm:enclose}) & \lean{HPC.Plane.enclose_holds}, from \lean{HPC.Plane.subset_component_of_far}, \lean{HPC.Plane.eastRay_subset_component}, \lean{HPC.Plane.compl_closedBall_subset_component}, \lean{HPC.Plane.exists_far_piece}, \lean{HPC.Plane.exists_uniform_pieces} and \lean{HPC.Plane.exists_mesh_compl_subset}.\\
Theorem~\ref{thm:polygon-around} & \lean{HPC.Plane.exists_polygon_around}; the splitting is \lean{HPC.Plane.exists_split_of_no_continuum}.\\[4pt]
\multicolumn{2}{@{}l}{\emph{Chapter~\ref{ch:tree}}}\\
Definition~\ref{def:treechain} & \lean{HPC.Plane.IsCoherent}, \lean{HPC.Plane.IsCircularChain}, \lean{HPC.Plane.IsTreeChain}, \lean{HPC.Plane.IsTreeLike}.\\
Lemma~\ref{lem:coherent-cover} & \lean{HPC.Plane.isCoherent_of_minimal_cover}.\\
Lemma~\ref{lem:filter-meet} & \lean{HPC.Plane.isTreeChain_filter_meet}.\\
Leaves (Lemma~\ref{lem:leaf}) & \lean{HPC.Plane.exists_leaf}, from \lean{HPC.Plane.ChainConn}, \lean{HPC.Plane.isCoherent_of_chainConn}, \lean{HPC.Plane.exists_chainConn_card}, \lean{HPC.Plane.circularChain_of_minimal} and \lean{HPC.Plane.no_two_nbrs}.\\
A logarithm on a fine tree chain (Theorem~\ref{thm:log-treechain}) & \lean{HPC.Plane.exists_log_on_treeChain}, by induction (\lean{HPC.Plane.exists_log_inv}) with \lean{HPC.Plane.exists_ball_branch}.\\
Theorem~\ref{thm:treelike-nonsep} & \lean{HPC.Plane.treeLike_nonseparating}.\\
Theorem~\ref{thm:log-treelike} & \lean{HPC.Hagopian.exists_log_of_isTreeLike}, from \lean{HPC.Hagopian.branchAt}, \lean{HPC.Hagopian.LogData}, \lean{HPC.Hagopian.LogData.branch_sub}, \lean{HPC.Hagopian.LogData.exp_const}, \lean{HPC.Hagopian.LogData.exists_glue} and \lean{HPC.Hagopian.log_mul_near_one}.\\
Lemma~\ref{lem:no-interior} & \lean{HPC.Plane.interior_eq_empty_of_noSeparating}.\\
Bing's offset (Lemma~\ref{lem:offset}) & \lean{HPC.Plane.gridSq}, \lean{HPC.Plane.gridVertices}, \lean{HPC.Plane.exists_grid_offset}.\\
Lemma~\ref{lem:notched} & \lean{HPC.Plane.notched}, \lean{HPC.Plane.isOpen_notched}, \lean{HPC.Plane.diam_notched_le}, \lean{HPC.Plane.inter_gridSq_subset_notched}, \lean{HPC.Plane.notched_diag1}, \lean{HPC.Plane.notched_diag2}, \lean{HPC.Plane.notched_adj}, \lean{HPC.Plane.notched_no_triple}.\\
Corollary~\ref{cor:notched-cover} & \lean{HPC.Plane.exists_far}, \lean{HPC.Plane.exists_notched_cover_bdd}.\\
Lemma~\ref{lem:compfamily} & \lean{HPC.Plane.compFamily}, \lean{HPC.Plane.compFamily_props}, \lean{HPC.Plane.connectedComponentIn_eq_of_meet}, \lean{HPC.Plane.compFamily_no_triple}.\\
A theta polygon (Lemma~\ref{lem:theta-polygon}) & \lean{HPC.Plane.exists_theta_polygon} (proved as \lean{HPC.Plane.theta_polygon_of_three}). The general threading of a circular chain of $m$ connected open sets, \lean{HPC.Plane.exists_threaded}, is not needed here.\\
Bing's point (Lemma~\ref{lem:bing-point}) & \lean{HPC.Plane.exists_theta_point_open}. The variant with polygonal $h_1,h_2$, where the polygonal form of Janiszewski's theorem suffices, is \lean{HPC.Plane.PlaneFacts.exists_theta_point}.\\
Proposition~\ref{prop:no-circular} & \lean{HPC.Plane.compFamily_no_circular_open}.\\
Bing's Theorem~6 (Theorem~\ref{thm:bing6}) & \lean{HPC.Plane.theorem6_direct}, via \lean{HPC.Plane.treeChain_of_cover_open}. The formal development also proves the theorem as a corollary of Theorem~7 below (\lean{HPC.Plane.theorem6_1951}); the two statements are identical.\\
Corollary~\ref{cor:treelike-iff} & \lean{HPC.Plane.isTreeLike_iff_nonseparating}.\\
Theorem~\ref{thm:weak-invariance} & \lean{HPC.Hagopian.exists_separating_subcontinuum_of_homeomorph}, with \lean{HPC.Hagopian.no_log_of_separates} for the point $p$.\\
Bing's Theorem~7 (Theorem~\ref{thm:bing7}) & \lean{HPC.Plane.theorem7_1951}; with the hypothesis ``$M$ tree-like'', as Bing states it, \lean{HPC.Plane.theorem7_1951_treeLike} (by Theorem~\ref{thm:treelike-nonsep}). A weak form with links that are open, connected, bounded, regular and non-separating is \lean{HPC.Plane.theorem7_weak}.\\
\end{longtable}
